\begin{qfdefinition}[Order-flow imbalance: Cont, Kukanov and Stoikov's event-based measure at the best quotes, against signed trade volume, and the family of imbalance features]\label{def:order-flow-imbalance}
\qffield{Statement} For consecutive best-quote states $(P^B_n,q^B_n,P^A_n,q^A_n)$, Cont, Kukanov and Stoikov define the contribution of the $n$th event to supply and demand at the best quotes,
\[ e_n=\mathbf 1\{P^B_n\ge P^B_{n-1}\}q^B_n-\mathbf 1\{P^B_n\le P^B_{n-1}\}q^B_{n-1}-\mathbf 1\{P^A_n\le P^A_{n-1}\}q^A_n+\mathbf 1\{P^A_n\ge P^A_{n-1}\}q^A_{n-1}, \]
so that at an unchanged bid $e_n$ is the change in bid size, a higher bid contributes the new bid size and a lower bid minus the old one, and the ask side enters with opposite signs; the \emph{order-flow imbalance} over $[t_{k-1},t_k]$ is $\mathrm{OFI}_k=\sum_{n=N(t_{k-1})+1}^{N(t_k)}e_n$, which ``encompasses trades, limit orders and cancelations'' \citep[arXiv v3, sec.~2.1, p.~4]{catalogue-0503}. The \emph{trade imbalance} is $\mathrm{TI}_k=\sum(b_n-s_n)$, buyer- minus seller-initiated volume (sec.~4.1, p.~16).

\qffield{Conditions} \emph{Evidence.} On one month (April 2010) of consolidated NYSE TAQ quotes and trades for 50 S\&P~500 stocks drawn at random, with 10-second intervals and regressions per half-hour, $\Delta P_k=\beta\,\mathrm{OFI}_k+\epsilon_k$ is linear with an average $R^2$ of 65 percent, against 32 percent for $\mathrm{TI}_k$; with both, the trade imbalance is significant in only 31 percent of subsamples (pp.~8 and~16), and the slope is inversely proportional to depth (abstract, p.~1). Timestamps were rounded to the second and quotes aggregated into an NBBO; Level~1 data only (p.~8). \emph{Names.} Bouchaud et al.\ call signed market-order volume $\sum\varepsilon_n\upsilon_n$ the ``order-flow imbalance'' \citep[chap.~11, sec.~11.4, p.~222]{catalogue-0528}; that is Cont et al.'s trade imbalance. A paper's ``OFI'' must be read from its formula.

\qffield{Practice} On a toy sequence, a market sell of 200 that is replenished, a cancellation of 300 at the ask and a market buy of 100 that empties the ask give $\mathrm{OFI}=+400$ and $\mathrm{TI}=-100$; the two can differ in sign. Trade imbalance, best-quote imbalance as in \hyperref[def:queue-imbalance]{queue imbalance}, multi-level depth imbalance, event imbalance by queue, and pressure aggregated across venues are a menu of constructions, not results.

\emph{Multi-level and cross-asset OFI}. Cont, Cucuringu and Zhang compute, for each of the top $M=10$ levels $m$, a per-event flow $\mathrm{OF}^{m,b}_n$ equal to $q^{m,b}_n$ if the level-$m$ bid price rose, $q^{m,b}_n-q^{m,b}_{n-1}$ if unchanged and $-q^{m,b}_n$ if it fell (and symmetrically on the ask), sum $\mathrm{OF}^{m,b}-\mathrm{OF}^{m,a}$ over each minute, and scale by the average depth of the first ten levels; the \emph{integrated} OFI is the first principal component of the ten scaled OFIs, with weights normalised to sum to one, which explains on average 89.06 percent of their variance (secs.~2.1--2.2, pp.~5--8, Table~2) \citep[arXiv v4, 13 June 2023]{wide-cont-cross-impact-ofi}.

Note that when a price moves away the printed per-level formula uses the \emph{new} size where Cont, Kukanov and Stoikov use the old one. Sample: Nasdaq ITCH data from LOBSTER for the 100 largest S\&P~500 stocks, 2017--2019, one-minute bins, 10:00 to 15:30, regressions per non-overlapping 30-minute window and out-of-sample $R^2$ on the next 30 minutes (p.~11). The universe is ``the top 100 components of S\&P 500 index, existing from 2017-01-01 to 2019-12-31'', selected ``based on their market capitalization as of the most recent market close at the time of our analysis, i.e.\ 2019-12-31'' (sec.~2.1 and n.~3, p.~5).

Contemporaneously, the best-level OFI explains 71.16 percent of returns in sample and 64.64 out of sample, the integrated OFI 87.14 and 83.83, and adding other stocks' OFIs (by LASSO) adds little: 73.87 and 66.03 for best-level, 87.85 and 83.62 for integrated (Tables~3 and~5, pp.~12 and~16). For forecasting one-minute-ahead returns, re-estimated on a rolling 30 minutes, every model has a \emph{negative} out-of-sample $R^2$ ($-0.37$ for own best-level OFIs, $-0.10$ with cross-asset terms), though forecast-implied portfolios, before trading costs, earn more with cross-asset terms, and the advantage decays within minutes (Table~8, p.~23; Table~11 and Fig.~10, pp.~26--27); integrated OFIs do not beat best-level ones for forecasting (p.~28). So ``multi-level'' OFI improves contemporaneous fit, not one-minute forecasts.

\emph{Universe selection}. Selecting by market capitalisation at the \emph{end} of the sample means that a stock's membership in 2017 and 2018 is decided by information from December 2019: the universe contains only firms that were large, listed and in the index at the end of 2019, and leaves out those that shrank, were acquired or left the index earlier. That is a look-ahead in the choice of stocks, separate from the temporal design, in which each regression is fitted on past minutes and evaluated on later ones (the next 30 minutes contemporaneously; rolling one-minute-ahead forecasts from the previous 30 minutes, pp.~11 and~22). The rolling design keeps parameters out of sample; it does not make the universe point-in-time.

The reported calculations are not thereby invalid, but the forecasting and portfolio results (most of all for 2017--2018, and the cross-sectional breadth of the portfolios) describe a survivor-selected set of large stocks, and their transfer to a universe formed with information available at each date is untested here; the paper reports no robustness check with a point-in-time universe. Normalising by volume, average daily volume or depth before comparing stocks is interpretation; Cont et al.'s depth result suggests depth.

\qffield{Pitfalls} Defining order-flow imbalance solely as signed trade volume conflates two measures: that is trade imbalance, not the event-based OFI. Cont et al.\ report contemporaneous fits, not forecasts, so they show neither that imbalance predicts the next quote move nor that it fails over longer horizons; on a subsample of stocks the results held for intervals from 10 quote updates to 10 minutes, the fit improving with the interval (p.~8). Any OFI feature used to forecast must be computed only from events before the forecast time.

\emph{Horizon and transfer in deep order-book models}. The LOBCAST benchmark re-ran 15 deep models on FI-2010 and on two LOBSTER sets it built, LOB-2021 (1--15 July 2021) and LOB-2022 (1--15 February 2022), each ten trading days of six NASDAQ stocks (SOFI, NFLX, CSCO, WING, SHLS, LSTR) chosen from the more liquid half of NASDAQ, ten levels sampled every ten events, split 6--2--2 days into training, validation and test, with up/stable/down labels from the mean mid-price over the next $k\in\{1,2,3,5,10\}$ observations against a fixed threshold \citep[arXiv v2, secs.~3--4.1, pp.~3--5]{w2-lobcast}.

F1 scores on the LOBSTER sets were 48--61 percent against up to 82.6 (BINCTABL) on FI-2010, BINCTABL losing about 19.6 percent on average across horizons; about half of hyperparameter runs diverged; and the authors' trading simulation (in a supplement) found profitability ``far from guaranteed'' (sec.~4.4, pp.~6--9). The horizon effect is not monotone: BINCTABL's correct-class rates on LOB-2021 were 54.5/65.5/57.0 percent (down/stable/up) at $k=1$ and 48.2/83.9/53.2 at $k=10$, the stable class gaining as the label smooths (Figure~2, p.~8).

Briola, Bartolucci and Aste label the mid-price change over 10, 50 and 100 book updates against one tick, on 15 large- and mega-cap NASDAQ stocks with LOBSTER data for 2017--2019, training on 45 consecutive days and testing on the next 10 each year, with five validation days drawn at random \emph{from the training period} \citep[arXiv v4, secs.~4--5, Table~2 and eq.~(2), pp.~7--8]{w2-briola-lob-guide}; they group stocks by mean spread (small tick $\gtrsim3$ ticks, large tick $\lesssim1.5$) and find predictability rising from small- to large-tick stocks (sec.~6). Two cautions follow (interpretation). Validation days inside the training window make model selection in-period rather than forward, and choosing stocks that stayed large- or mega-cap through 2019 conditions on the period's outcome, as with Cont, Cucuringu and Zhang's universe above.

Their transaction metric $p_T=CT/(PT+TT-CT)$, the overlap between round trips implied by predicted and by realised labels, is described by the authors as ``assumption-free'': it assumes no costs, fills or latency (sec.~7, eq.~(3), pp.~22--23), so it measures whether entries and exits are placed correctly, not ``profitable round trips after spread''; small-tick stocks at the 10-update horizon gave $p_T$ of about 0.06 (CHTR, GOOG, GS) and 0.12 (IBM, MCD, NVDA) without probability thresholds (p.~24).

\emph{Order-flow inputs across horizons} (SSRN version of 6 August 2021, not the 2023 journal version). Kolm, Turiel and Westray train one network per stock on LOBSTER data for 115 Nasdaq stocks, January 2019 to January 2020, chosen as Nasdaq-100 members at some time in the period with at least 260 trading days and no corporate actions over it (so, again, a selection that uses the whole period), defining ten-level book states (LOB), order flows (OF) and OFI and comparing six models trained on LOB or on OF inputs \citep[SSRN version, secs.~2.1 and~3.1, PDF pp.~5--7 and~12]{wide-kolm-deep-ofi}.

Horizons are stock-specific: $\Delta t=2.34\cdot10^7/N$ milliseconds, with $N$ the average daily number of non-zero mid-price changes, $h_k=k\Delta t/5$ for $k=1,\dots,10$, and returns measured after a ten-millisecond latency buffer; windows of one validation week, four training weeks and one test week roll forward by three weeks, 18 windows in all (sec.~3.2, PDF pp.~14--15). In pairwise comparisons of each model trained on the two inputs, a Table~3 entry is, for one model and horizon, the number of stocks with a significantly positive $t$-statistic (1 percent level) minus the number with a significantly negative one (sec.~4.1.2, PDF p.~19; caption, PDF p.~21).

The CNN-LSTM row runs from 70 at 0.2 to 93 at 1.8 average price changes (89 on the horizon average), the LSTM rows 114--115. A net count is not a count of winning stocks, and $115$ minus it is not a count of losing ones: stocks with no significant difference count on neither side, and the raw counts are not reported. The source is ambiguous here, and no exact count is stated in this entry: the caption and text also read an entry $X$ as OF outperforming ``for $X$ stocks'', the caption calls the tested quantity a ``loss difference'' where the text says a difference in $R^2_{OS}$, and the text's ``about 40 stocks for which the CNN-LSTM with LOB inputs perform better'' at the shorter horizons (PDF p.~20) cannot be derived from the net counts.

Out-of-sample $R^2$ rises up to about two average price changes (Fig.~2, PDF p.~18, whose horizons stop at two, so it cannot show what follows); the decline beyond two is shown by Fig.~8 (PDF p.~28), on horizons to ten: read from the plot, CNN-LSTM peaks near 1.5 percent at about two and falls to roughly 0.64 percent at ten, although the text describing it says $R^2_{OS}$ ``plateaus at around two average price changes'' (sec.~4.3, PDF p.~27). These are statistical forecasts with no costs; the authors say that for trading strategies ``proper modeling of transaction costs and application of an appropriate execution strategy is critical to their profitability'' (PDF p.~31). None of this establishes that any signal here pays after costs.

\emph{Judging a signal by what it is worth after trading}. Briola, Bartolucci and Aste list the assumptions that a history-only backtest of a deep-learning forecast must make: (i) the infrastructure to record and process live data, forecast and execute in time; (ii) always being executed; (iii) zero market impact; (iv) zero transaction costs, whose combination ``ruins any attempt to produce a reliable backtest'' \citep[arXiv v4, sec.~7, pp.~22--23]{w2-briola-lob-guide}; they also report how often their tick-time horizons (10, 50, 100 updates) take under one second, one to ten seconds or longer in clock time, because practitioners care about ``actual practicability in real-world scenarios'' (Table~7, pp.~16--17).

LOBCAST's own trading simulation found profitability ``far from guaranteed'' (above), and Kolm et al.\ make costs and execution ``critical to their profitability'' (above); Gould and Bonart's classification scores (\hyperref[def:queue-imbalance]{queue imbalance}) and Moallemi and Sağlam's latency cost give the predictability and the delay penalty but no net result. So the checklist's items on tradability (spread, fees, fill and queue loss), latency and tick size now have sources; its capacity and crowding items do not.

\emph{A costed test, specified but not run.} Data: message and book files (for example LOBSTER, ten levels) for named stocks and days, under the provider's licence; LOBSTER's free samples require a request and acceptance of terms. Signal: queue imbalance or event-based OFI computed only from events before the decision time. Decision: trade when the signal passes a threshold chosen on training days; order: a market order at the touch after a stated latency (in events and in milliseconds), filled up to the displayed size, the rest unfilled; exit: at the next mid-price change or a fixed horizon, again crossing the spread. Costs: the spread paid on entry and exit and the venue's taker fee at the data date.

Validation: forward by day, with each threshold and every variant logged, reporting net profit per trade and per day, the Sharpe ratio with \hyperref[meth:sharpe-ratio-standard-error]{Lo's standard error} and the \hyperref[meth:deflated-sharpe-ratio]{deflated Sharpe ratio} over the variants tried, and the share of displayed depth used. Until such a test is run on licensed data, no signal here is shown to pay after costs.

\qffield{See also} \hyperref[meth:trade-sign-classification]{trade-sign classification}; \hyperref[def:queue-imbalance]{queue imbalance}; \hyperref[meth:information-driven-bars]{information-driven bars} (imbalance bars, a different use); \hyperref[mod:kyle-single-auction]{Kyle} (impact linear in net order flow); \hyperref[meth:order-flow-impact-regression]{order-flow impact regression}.

\qffield{Further reading} Kolm, Turiel and Westray, SSRN working-paper version of 6 August 2021 (Pitfalls); later published in \emph{Mathematical Finance}.

\end{qfdefinition}
